\chapter{Chapter06 — Small Bias Fourier Estimator}
%%% generated by the blueprint pipeline from FABL.Chapter06.LearningAndTesting.SmallBiasFourierEstimator
\section{Overview}

Declarations of the Lean module \texttt{FABL.Chapter06.LearningAndTesting.SmallBiasFourierEstimator}.

\section{Declarations}

\begin{definition}[Empirical Fourier estimate]
\lean{FABL.smallBiasFourierEstimate}
\statementsource{Analysis-of-Boolean-Functions-by-Ryan-ODonnell}{
  pdf-ffbe9c7b3306-p157-b008,
  pdf-ffbe9c7b3306-p157-b009
}
\label{FABL.smallBiasFourierEstimate}
\leanok
\uses{FABL.F₂Cube, FABL.binaryCubeSignEquiv, FABL.monomial}
Let $f$ be a real-valued function on the $n$-dimensional sign cube, let $x_1,\ldots,x_m$
be points of the $n$-dimensional Boolean cube, and let $U\subseteq\{1,\ldots,n\}$. The
empirical Fourier estimate is $\frac{1}{m}\sum_{i=1}^m f(\chi(x_i))\,\chi_U(\chi(x_i))$,
where $\chi$ is the coordinatewise sign map and $\chi_U$ is the character indexed by
$U$.
\end{definition}

\begin{definition}[Estimator work bound]
\lean{FABL.smallBiasFourierEstimatorWork}
\statementsource{Analysis-of-Boolean-Functions-by-Ryan-ODonnell}{
  pdf-ffbe9c7b3306-p157-b008,
  pdf-ffbe9c7b3306-p157-b009
}
\label{FABL.smallBiasFourierEstimatorWork}
\leanok
For a natural number sample size $m$ and a set $U$ of coordinates, the estimator's work
bound is $m(|U|+2)+1$.
\end{definition}

\begin{definition}[Estimator cost]
\lean{FABL.smallBiasFourierEstimatorCost}
\statementsource{Analysis-of-Boolean-Functions-by-Ryan-ODonnell}{
  pdf-ffbe9c7b3306-p157-b008,
  pdf-ffbe9c7b3306-p157-b009
}
\label{FABL.smallBiasFourierEstimatorCost}
\leanok
\uses{FABL.DeterministicQueryProgram.queryBatchCost, FABL.LearningCost, FABL.smallBiasFourierEstimatorWork}
For a natural number sample size $m$ and a set $U$ of coordinates, the estimator cost is
the sum of the cost of a batch of $m$ deterministic membership queries and a cost with
zero random-example calls, zero membership queries, and $m(|U|+2)+1$ local computation
steps.
\end{definition}

\begin{definition}[Fourier estimator query program]
\lean{FABL.smallBiasFourierEstimatorProgram}
\statementsource{Analysis-of-Boolean-Functions-by-Ryan-ODonnell}{
  pdf-ffbe9c7b3306-p157-b008,
  pdf-ffbe9c7b3306-p157-b009
}
\label{FABL.smallBiasFourierEstimatorProgram}
\leanok
\uses{FABL.DeterministicQueryProgram, FABL.F₂Cube, FABL.binaryCubeSignEquiv, FABL.monomial, FABL.smallBiasFourierEstimatorWork}
Let $x_1,\ldots,x_m$ be points of the $n$-dimensional Boolean cube and let $U$ be a
subset of the $n$ coordinates. The program queries the sign encodings of all $m$ points,
then uses the answers $y_1,\ldots,y_m$ to return $\frac{1}{m}\sum_{i=1}^m
y_i\chi_U(\chi(x_i))$, charging $m(|U|+2)+1$ local computation steps; here $\chi$ is the
coordinatewise sign map and $\chi_U$ is the character indexed by $U$.
\end{definition}

\begin{theorem}[Cost of the Fourier estimator program]
\lean{FABL.DeterministicQueryProgram.runWithCost_smallBiasFourierEstimatorProgram}
\label{FABL.DeterministicQueryProgram.runWithCost_smallBiasFourierEstimatorProgram}
\leanok
\uses{FABL.DeterministicQueryProgram.runWithCost, FABL.F₂Cube, FABL.smallBiasFourierEstimate, FABL.smallBiasFourierEstimatorCost, FABL.smallBiasFourierEstimatorProgram}
\difficulty{1}
Let $m$ be a natural number, let $f$ be a real-valued function on the $n$-dimensional
sign cube, let the sample be a map from an $m$-element index set to binary vectors of
length $n$, and let $U$ be a finite set of coordinates. Running the small-bias Fourier
estimator program with oracle $f$ returns the sample average of $f(x)$ times the
monomial indexed by $U$, where each binary sample is converted to a sign vector,
together with its prescribed cost: zero random examples, $m$ membership queries, and
$m(|U|+2)+1$ units of local work.
\end{theorem}

\begin{theorem}[Multiplication by a monomial shifts Fourier support]
\lean{FABL.fourierCoeff_pointwise_mul_monomial}
\label{FABL.fourierCoeff_pointwise_mul_monomial}
\leanok
\uses{FABL.fourierCoeff, FABL.monomial, FABL.monomial_mul_monomial}
\difficulty{3}
Let $n$ be a nonnegative integer, let $f:\{-1,1\}^n\to\mathbb R$, and let
$U,S\subseteq\{1,\ldots,n\}$. The Fourier coefficient at $S$ of the function
$x\mapsto f(x)\prod_{i\in U}x_i$ equals $\widehat f(U\triangle S)$.
\end{theorem}

\begin{theorem}[Invariance of the Fourier one-norm]
\lean{FABL.fourierOneNorm_pointwise_mul_monomial}
\label{FABL.fourierOneNorm_pointwise_mul_monomial}
\leanok
\uses{FABL.fourierCoeff_pointwise_mul_monomial, FABL.fourierOneNorm, FABL.monomial}
\difficulty{4}
For a natural number $n$, a real-valued function $f$ on the $n$-dimensional sign cube,
and a set $U$ of coordinates, the Fourier one-norm of the pointwise product of $f$ with
the monomial indexed by $U$ equals the Fourier one-norm of $f$.
\end{theorem}

\begin{theorem}[Mean and Fourier coefficient]
\lean{FABL.mean_pointwise_mul_monomial}
\proofsource{Analysis-of-Boolean-Functions-by-Ryan-ODonnell}{pdf-ffbe9c7b3306-p157-b009}
\label{FABL.mean_pointwise_mul_monomial}
\leanok
\uses{FABL.fourierCoeff, FABL.fourierCoeff_pointwise_mul_monomial, FABL.mean, FABL.mean_eq_fourierCoeff_empty, FABL.monomial}
\difficulty{3}
Let $f$ be a real-valued function on the $n$-dimensional sign cube, and let $U$ be a set
of coordinates. The mean of the pointwise product of $f$ with the monomial that
multiplies the signs in $U$ equals the Fourier coefficient of $f$ indexed by $U$.
\end{theorem}

\begin{theorem}[Empirical Fourier estimate as an expectation]
\lean{FABL.smallBiasFourierEstimate_eq_uniformPushforward_expectation}
\label{FABL.smallBiasFourierEstimate_eq_uniformPushforward_expectation}
\leanok
\uses{FABL.F₂Cube, FABL.ProbabilityDensity.expectation, FABL.ProbabilityDensity.expectation_uniformPushforward, FABL.ProbabilityDensity.uniformPushforward, FABL.binaryCubeSignEquiv, FABL.monomial, FABL.smallBiasFourierEstimate}
\difficulty{3}
Let $n$ be a natural number, let $m$ be a positive natural number, let $f$ be a
real-valued function on the sign cube of dimension $n$, let the sample consist of $m$
points of the binary cube of dimension $n$, and let $U$ be a set of coordinates. The
empirical average of $f$ times the sign monomial indexed by $U$ over the encoded sample
equals the expectation of that same function under the density obtained by pushing
forward the uniform distribution on the sample points.
\end{theorem}

\begin{theorem}[Small-bias Fourier estimation error]
\lean{FABL.abs_smallBiasFourierEstimate_sub_fourierCoeff_le}
\proofsource{Analysis-of-Boolean-Functions-by-Ryan-ODonnell}{
  pdf-ffbe9c7b3306-p157-b003,
  pdf-ffbe9c7b3306-p157-b004,
  pdf-ffbe9c7b3306-p157-b009
}
\label{FABL.abs_smallBiasFourierEstimate_sub_fourierCoeff_le}
\leanok
\uses{FABL.F₂Cube, FABL.ProbabilityDensity.IsBiased, FABL.ProbabilityDensity.abs_expectation_signFunction_sub_mean_le_fourierOneNorm_mul, FABL.ProbabilityDensity.uniformPushforward, FABL.fourierCoeff, FABL.fourierOneNorm, FABL.fourierOneNorm_pointwise_mul_monomial, FABL.mean_pointwise_mul_monomial, FABL.monomial, FABL.smallBiasFourierEstimate, FABL.smallBiasFourierEstimate_eq_uniformPushforward_expectation}
\difficulty{3}
Let $m$ be a positive integer, let $f:\{\pm1\}^n\to\mathbb R$, let
$U\subseteq\{1,\ldots,n\}$, and let a sample of $m$ points in the binary cube induce a
probability density whose nonconstant Fourier coefficients have absolute value at most
$\delta$, where $\delta\ge0$. Then the absolute difference between the sample estimate
of the Fourier coefficient of $f$ at $U$ and its uniform Fourier coefficient is at most
$\delta$ times the sum of the absolute values of all Fourier coefficients of $f$.
\end{theorem}

\begin{theorem}[Parameterized small-bias Fourier estimation error]
\lean{FABL.abs_smallBiasFourierEstimate_sub_fourierCoeff_le_parameter}
\proofsource{Analysis-of-Boolean-Functions-by-Ryan-ODonnell}{
  pdf-ffbe9c7b3306-p157-b003,
  pdf-ffbe9c7b3306-p157-b004
}
\label{FABL.abs_smallBiasFourierEstimate_sub_fourierCoeff_le_parameter}
\leanok
\uses{FABL.F₂Cube, FABL.ProbabilityDensity.IsBiased, FABL.ProbabilityDensity.uniformPushforward, FABL.abs_smallBiasFourierEstimate_sub_fourierCoeff_le, FABL.fourierCoeff, FABL.fourierOneNorm, FABL.smallBiasFourierEstimate}
\difficulty{3}
Let $m$ be a positive integer, let $f:\{\pm1\}^n\to\mathbb R$, let $U$ be a subset of
the $n$ coordinates, and let $\varepsilon,s\in\mathbb R$ satisfy
$0<\varepsilon\le\tfrac12$ and $1\le s$. Suppose the sample's uniform pushforward
density has every nonzero binary Fourier coefficient bounded in absolute value by
$\varepsilon/s$, and the sum of the absolute values of the Fourier coefficients of $f$
is at most $s$. Then the absolute difference between the sample estimate of the Fourier
coefficient of $f$ at $U$ and its uniform Fourier coefficient is at most $\varepsilon$.
\end{theorem}
